\subsection{Rainbows}\label{subsec:rainbows}

Of course light reflects off of everything, not just flat mirrors.  A
general discussion of that phenomenon is more properly done in a
course on classical physics.  We will limit ourselves to the geometry
involved when a light beam strikes a spherical mirror, like the
surface of a raindrop. Since a raindrop  is not a perfect mirror some
of the beam passes through the mirror and refracts, while some of it
reflects back off of the surface.

\begin{wrapfigure}[]{r}{2in}
  \captionsetup{labelformat=empty}
  \centerline{\includegraphics*[height=1in,width=2in]{../Figures/Refraction10}}
  \label{fig:Refraction10}
\end{wrapfigure}
Happily nothing changes much. We need only observe that light only
interacts with the mirror locally, where it actually strikes  the
mirror. And since we know that at very small scales any curve is
indistinguishable from it's tangent line it should be clear that when
a light beam strikes a curved mirror it reflects or refracts exactly
as if it were striking a flat mirror tangent to the curve at the point
of interaction as seen at the right. This is again the Principle of Local
Linearity  

\begin{wrapfigure}[]{l}{2.1in}
  \captionsetup{labelformat=empty}
  \centerline{\includegraphics*[height=.75in,width=2.1in]{../Figures/DoubleRainbow6}}
  \label{fig:DoubleRainbow6}
\end{wrapfigure}
The ordering of the colors of a rainbow, the shape of a
rainbow, its position relative to the viewer, and the simple fact that
rainbows exist to be seen are a result of the reflective and
refractive properties of light. Even better for our purposes
optimization is needed to answer them.

Newton determined experimentally that white light was a mixture of
other colors. It was well known that the time that a beam of light
passing through a prism would emerge with different colors bent
(refracted) at different angles as seen above on the left. A prevailing
theory of the time was that the prism had altered the light in some
way that had added the colors. Thus it was believed that the color
bands that emerged from the prism were a property of the prism, not of
the light.

\begin{wrapfigure}[]{r}{4in}
  \captionsetup{labelformat=empty}
  \label{fig:DoubleRainbow9}
  \centerline{\includegraphics*[height=1.9in,width=4in]{../Figures/DoubleRainbow9}}
  \caption{The double rainbow on the left was captured by Elizabeth
    Eschel (a friend of one of the authors) of Tsfat/Safed, Israel in
    2018 while on her way to
    work.}\end{wrapfigure}
Newton debunked that theory by passing \emph{only} the red band
through a second prism. He reasoned that if the first prism had
somehow altered the light then the second would alter it again. But
when only red light emerged from the second prism, refracted but
otherwise unchanged, he concluded that the color was an inherent
property of the light, not the prism.

% wrote a seminal work on optics and .  This can be seen as light passes
% through a prism or in a rainbow.  Surely you heard the mnemonic
% \alert{ROY G BIV} for the colors of a rainbow.  Which color is on top,
% red or violet?  Before you go further, take a guess.  % \\
Look closely at the three rainbows in the images at the right above. They
share certain properties that we would like to understand. These are:
\begin{enumerate}
\item They are bright. At least they are brighter than the surrounding
  viewscape. Indeed it is only the brightness of the light coming from
  the apparent location of the rainbow that makes it visible since a
  rainbow is not a physical object.
\item They have a particular shape, but what is it? Are they
  circular? Parabolic? Elliptical? Some other, possibly unnamed
  shape?
\item Two of the rainbows shown are double rainbows. Notice that the
  secondary rainbow is dimmer\footnote{That's actually what makes it
    the secondary.}. Also the colors are reversed from the primary to
  the secondary. We'd like to understand that.
\end{enumerate}
Each of the properties listed above can be couched as an optimization
problem. We will address them one at a time.

To keep things as simple as possible we will assume that water
droplets suspended in the atmosphere assume a perfectly spherical
shape. This is close enough to reality that the errors in our
computations will not be significant.


Also, we will analyze this as if everything were two-dimensional. That
is we will draw the light beam crossing into a circular raindrop
rather than a spherical raindrop. Again this is just to keep things
simple. The analysis of the three dimensional problem  yields
similar results but is much harder algebraically.

\begin{wrapfigure}[]{r}{1.7in}
  \captionsetup{labelformat=empty}
  \centerline{\includegraphics*[height=1.7in,width=1.7in]{../Figures/Refraction6}}
  \label{fig:Refraction6}
\end{wrapfigure}
As we've seen when light refracts across the plane dividing air and
water it bends toward the normal (perpendicular). The physical reason
for this is that the velocity of light in air is faster than its
velocity in water. If it transitions from a slower medium (like water)
to a faster medium (like air) then it bends away from the normal. Thus
in the diagram at the right we can deduce that the velocity is lower
in the region below the boundary (in black) whichever direction the
light ray (in red) is moving.



% Of course the sketch above doesn't capture the curvature of the rain
% droplet.  Our sketch should look more like the figure at the left,
% where the light beam is crossing the black, circular boundary of the
% droplet. Does Snell's Law still apply when the boundary is curved?

% Of course it does. The light beam doesn't care what shape the boundary
% has away from the point of entry. It only interacts with the boundary
% locally, and locally any curve is indistiquishable from it's tangent
% line so the light beam will refract as it crosses into the droplet
% just as if it were crossing the line tangent to the circle at the
% point of intersection. As we saw in\notetoself{We need to add this
% problem somewhere before this, if we don't already have it: Show that the line tangent to a
% circle is orthogonal to the radius of the circle.  -- Sep 3, 2020}
% Problem~\#\ref{} the line tangent to a circle at a point is orthogonal
% to the radius that passes through the same point. The radius and
% tangent are the blue, dotted lines in the diagram at the left.
But the physics involved is not our concern. All we need is  the
knowledge that Snell's Law is in play in the diagram at the left.
$\frac{\sin(\alpha)}{v_1}=\frac{\sin(\beta)}{v_2}.$
The green line in the diagram is the path the light beam would have
followed had it \emph{not} refracted at the boundary.

\begin{ProficiencyDrill}
  For later reference notice that the angle between the red and green
  lines is $\alpha-\beta$ as indicated. \hint{This is an elementary
    result from geometry. Don't make it hard.}
\end{ProficiencyDrill}
We now have everything we need to trace the path of a beam of light
into and back out of a water droplet.
In the sketch below the red arrows indicate  the path the light
beam actually follows, $O$ is
the center of the droplet, and the line segments $OA,
OB,$ and $OC$ are all radii of the droplet.

% \begin{figure}[h]
  
  \centerline{\includegraphics*[height=2.5in,width=4in]{../Figures/Refraction7}}
  
%   \caption{The path of a typical light beam as it moves through a
%     spherical droplet.}
%   \label{fig:DoublePath}
% \end{figure}
 
In order to address the question of the rainbow's brightness -- why it
is even visible in the first place -- we want to find the total
clockwise rotation of the direction of the beam from the time it
enters the raindrop to the time it exits. 

\begin{wrapfigure}[]{r}{1.5in}
  \captionsetup{labelformat=empty}
  \centerline{\includegraphics*[height=.6in,width=1.5in]{../Figures/Refraction8}}
  \label{fig:Refraction8}
\end{wrapfigure}
That is, if we place the
starting and ending beams head-to-tail as indicated in the sketch at the right we
want to find the angle $D(\alpha)$. This is called the
\term{deviation angle} and for a given entry angle $\alpha$ it tells
us what the exit angle will be. We are interested in finding the value
of $\alpha$ that \emph{minimizes} $D(\alpha)$.

Here's why. Suppose $D(\alpha)$ is minimum for some angle, say
$\alpha=\phi$, and that $D(\phi)=\psi$ is the minimum. Then all of the
light beams entering the droplet at an angle \emph{near} $\phi$ will
exit \emph{near} $\psi$. It is generally true that things which start
out close together usually end up close together. Functions with this
property are called \term{continuous functions}.

\begin{wrapfigure}[]{l}{3.5in}
  \captionsetup{labelformat=empty}
  \centerline{\includegraphics*[height=1.5in,width=3.5in]{../Figures/DoubleRainbow3}}
  \label{fig:DoubleRainbow3}
\end{wrapfigure}
But the nature of the function
$D(\alpha)$ guarantees that there is a large interval of entry
angles which will all have approximately the same exit angle.  The
easiest way to see this is to look at the the shape of the graph of
$D(\alpha)$ below. Notice that the
black box we drew around then minimum point (denoted\aside{For now you
  will have to trust that we have represented $D(\alpha)$
  accurately. We will shortly be deriving a formula for $D(\alpha)$
  and you can verify the accuracy of this graph yourself.}
$(\gamma,\psi)$ in the sketch.). is much wider than it is
tall.
% This says that all of the the light beams entering at an angle
% between about $0.8$ radians and $1.2$ radians will exit at about $2.4$
% radians. That is the light will be concentrated (brighter) at a
% deviation angle of about $2.4$ radians (about $138^\circ$. 

% We will soon see that the deviation angle
% depends \emph{only} on the angle $\alpha$ as we've
% indicated. 
From the diagram it is clear that for entry angles between a little
more than $0.8$ radians ($\approx45.8^\circ$) and a little more than
$1.2$ radians ($\approx68.8^\circ$) the exit angles will all be very
close to $2.4\approx138^\circ$ radians. Outside of the interval
$(.8, 1.2)$ the exit angles grow very quickly.

This says that the light beams coming out of a droplet at an angle of
$\psi\approx2.4$ radians are a concentration of the beams that entered
between about $.8$ and $1.2$ radians. So light
exiting at about $2.4$ radians will be brighter than the surrounding
light. This is what makes a rainbow visible in the first place.

Now of course, we need to do the calculations to confirm all of
this. To find a formula for the deviation angle, $D(\alpha)$ we
follow the path of a typical light beam as it refracts and reflects
its way into and then back out of a typical droplet.

In Figure~\ref{fig:DoublePath} we see that the beam (in red) makes a
turn of $\alpha-\beta$ at the point $A$ and another turn of the
same magnitude at point $C$ so
$$
D(\alpha) =
2(\alpha-\beta) + \{\text{the magnitude of the rotation at point
  $B$}\}.
$$

\begin{wrapfigure}[]{r}{2in}
  \captionsetup{labelformat=empty}
  \centerline{\includegraphics*[height=1.8in,width=2in]{../Figures/Refraction9}}
  \label{fig:Refraction9}
\end{wrapfigure}
To find the magnitude of the rotation at point $B$ consider the sketch at
the right. Clearly at point $B$ the beam turns by the angle labeled
$\theta$.

\begin{ProficiencyDrill}
  \begin{enumerate}[label={\bf{}(\alph*)}]
  \item   Use Geometry to show that $\alpha=\angle AOB$.
  \item Show that $\alpha=\pi-2\beta$. \hint{The interior angles
      of triangle $OAB$ must sum to $\pi$. What is $\beta+\theta$
      equal to?}
  \end{enumerate}
\end{ProficiencyDrill}
Therefore
\begin{align*}
  D(\alpha) &= 2(\alpha-\beta) + \{\text{the magnitude of the
              rotation at point  $B$}\}\\
            &= 2(\alpha-\beta) + \alpha\\
            &= 2(\alpha-\beta) + \pi-2\beta, \text{ or}\\
            &= \pi+2\alpha-4\beta.
\end{align*}

Notice that $\alpha$ and $\beta$ are not independent, they are governed by Snell's
Law which we will rewrite as
\begin{equation}
  \sin(\alpha)=\frac{v_{\alpha}}{v_{\beta}}\sin(\beta)=k\sin(\beta).\label{eq:IndRef}
\end{equation}
The constant $k=\frac{v_{\alpha}}{v_{\beta}}$ is called  the ``index of refraction'' (for
water).

\begin{ProficiencyDrill}
  Show that $\beta=\inverse\ \lt{}sin\left(\frac{\sin(\alpha)}{k}\right)$.
\end{ProficiencyDrill}
Knowing $\beta$ we can now express the deviation angle,
$D(\alpha)$ as a function of $\alpha$ alone:
\begin{align*}
  D(\alpha) &= \pi+2\alpha-4\beta \\
  D(\alpha) &= \pi+2\alpha-4\inverse\sin\left(\frac{\sin(\alpha)}{k}\right).
\end{align*}

\begin{embeddedproblem}
  For visible light  the index of refraction
  is $k\approx\frac43$. Mathematically, the problem becomes minimizing
  $D(\alpha)=\pi+2\alpha-4\beta$ subject to the constraint
  $\sin(\alpha)=k\sin(\beta)$, $ 0\le\alpha\lt{}\frac{\pi}{2} $.

  \begin{enumerate}[label={\bf{}(\alph*)}]
  \item Graph $D(\alpha)$ with $k=\frac43$ and confirm that the graph
    we showed you earlier is correct.
  \item Use Newton's Method to show that the coordinates,
    $(\gamma,\psi),$
    of the minimum point on the graph are $\gamma\approx1.04\approx59.59^\circ$ and
    $\psi\approx2.41\approx138.1^\circ$ respectively.
  \end{enumerate}
\end{embeddedproblem}
% ****************** Pick up here *************************************

% Graphing $D(\alpha)$ for this value of  $k$ and
% $0\le{}\alpha{}\lt{}\frac{\pi{}}{2}$, we have.
% Notice from the graph that the minimum of this curve appears to be
% around the point $(1.03, 2.41)$.  This says that the minimum value of
% $D(\alpha)$  is about $2.41\text{ radians }\approx 138^\circ$  and it
% occurs when $\alpha\approx1.03\text{ radians }\approx59.4^\circ$.
% As we can see from our sketch the  minimum 
% value of $\alpha$ near  $1.03$  the amount of deviation is roughly
% the same.

% \vskip.5in{}
% ***************************************************************\\
% As you can see in the diagram below, the ray of light entering the
% raindrop at an angle of radian measure $\alpha$ is refracted to an
% angle of radian measure $\beta$.  This is governed by Snell's Law.\\
% % \begin{wrapfigure}[]{l}{3in}
% %   \captionsetup{labelformat=empty}
% %   \centerline{\includegraphics*[height=1.75in,width=3.25in]{../Figures/DoubleRainbow2}}
% %   \label{fig:DoubleRainbow2}
% % \end{wrapfigure}
% $\frac{\sin(\alpha)}{v_\alpha}=\frac{\sin(\beta)}{v_\beta}$, where
% $v_\alpha$ and $v_\beta$ are the velocities of light in air and water,
% respectively.  Actually, only a portion is refracted as the other
% portion is reflected off the exterior of the raindrop.  Once inside
% the raindrop, it reflects off of the back of the raindrop (again only
% a portion) so that the angle of incidence equals the angle of
% reflection.  A portion of the light then exits the raindrop governed
% by Snell's Law again.  The total amount of (clockwise) rotation in
% this process is called the deviation angle, $D(\alpha)$.%

% \begin{embeddedproblem}
%   Show that $D(\alpha)=\pi-4\beta+2\alpha.$
% \end{embeddedproblem}


\begin{wrapfigure}[]{r}{2.5in}
  \captionsetup{labelformat=empty}
  \includegraphics*[height=.8in,width=2.5in]{../Figures/DoubleRainbow5}
  \label{fig:DoubleRainbow5}
\end{wrapfigure}
This says that light rays entering the raindrop at approximately
$59^\circ$ have the highest concentration when exiting the raindrop.
At other values of $\alpha$ the dispersion of light is greater.  This
provides the angle that the rainbow makes with an observer.  This
supplementary angle to $D(\alpha)$ is called the rainbow angle,
$$
R(\alpha) =\pi-D(\alpha)=\pi-(\pi+2\alpha-4\beta)=4\beta-2\alpha\text{
  radians.}
$$
For white light, we have a rainbow angle of
$$
R(1.03)\approx\pi-2.41\approx0.734\text{ radians}\approx42^\circ.
$$

To obtain the spread of colors, note that the index of refraction
changes with the color of light.  We said earlier that the index of
refraction for visible light is approximately $\frac43$. This is true,
but it actually varies with the frequency (color) of the light. The
red light refracts at one angle and the orange at a slightly different
angle so they appear in the rainbow adjacent to each other. This
variation is continuous throughout the spectrum which is what gives us
the rainbow.


\begin{embeddedproblem}
  We could use Newton's Method like we did earlier to find the
  minimum of $D(\alpha)$ for various values of $k$ but it is a
  little easier to do it this way:
     \begin{enumerate}[label={\bf{}(\alph*)}]
  \item Show that $D(\alpha)$ is minimized when % $\cos(\alpha)=\frac{k}{2}\cos(\beta)$.
  \item Use the equation obtained in part (a) along with the
    constraint $\sin(\alpha)=k\sin(\beta)$ to show that $D(\alpha)$
    is minimized when $\sin(\beta)=\sqrt{\frac{4-k^2}{3k^2}}$ and
    $\sin(\alpha)=\sqrt{\frac{4-k^2}{3}}$.  \hint{Solving for one
      variable in one equation and substituting into the other is
      very difficult. Instead try squaring both equations and add
      them together? From there solving for $\sin(\beta)$ is
      straightforward\footnote{Remember that ``straightforward''
        does not mean easy.}.}
  \item Use the results from part (b) to fill in the following
    table.
    \begin{table}[h]
      \centering
      \begin{tabular}{|c|c|c|c|c|c|c|} \hline
        Color& k& $\beta$ in degrees&$\alpha$ in degrees&$D(\alpha)$ in degrees&$R(\alpha)$ in degrees\\\hline
        Red & 1.331 & &&&\\\hline
        Orange & 1.332 & &&&\\\hline
        Yellow & 1.333 & &&&\\\hline
        Green & 1.335 & &&&\\\hline
        Blue & 1.337 & &&&\\\hline
        Indigo & 1.340 & &&&\\\hline
        Violet & 1.344 & &&&\\\hline
      \end{tabular}
    \end{table}
  \end{enumerate}
\end{embeddedproblem}

% \begin{wrapfigure}[]{r}{1.25in}
%   \captionsetup{labelformat=empty}
%   \centerline{\includegraphics*[height=.in,width=1.25in]{../Figures/DoubleRainbow7}}
%   \label{fig:DoubleRainbow7}
% \end{wrapfigure}
The values of $R(\alpha)$ you obtained in the table
accounts for the various color bands that appear in the rainbow as
sunlight hits droplets at various angles.% \\
% \centerline{\includegraphics*[height=1.in,width=3in]{../Figures/DoubleRainbow7}}

The primary rainbow is created by light entering the upper portion of
the raindrop.  If light is bright enough then a secondary rainbow is
created from light entering the lower portion of the raindrop as in
the diagram below.\\
\centerline{\includegraphics*[height=2in,width=3.3in]{../Figures/DoubleRainbow8}}
This time the deviation angle $D(\alpha)$ is the total amount of counterclockwise
rotation and the rainbow angle would be $R(\alpha)=D(\alpha)-\pi$.  Note that the extra
reflection accounts for the rainbow being not as bright.

\begin{embeddedproblem}
  Show that in the case of the secondary rainbow
  $$
  D(\alpha)=2\pi+2\alpha-6\beta
  $$
  and show that this is minimized when
  $\sin(\beta)=\sqrt{\frac{9-k^2}{8k^2}}$ and
  $\sin(\alpha)=\sqrt{\frac{9-k^2}{8}}$.
  Use this to fill in the following table:
  \begin{table}[h]
    \centering
    \begin{tabular}{|c|c|c|c|c|c|c|} \hline
      Color& k& $\beta$ in degrees&$\alpha$ in degrees&$D(\alpha)$ in degrees&$R(\alpha)$ in degrees\\\hline
      Red & 1.331 & &&&\\\hline
      Orange & 1.332 & &&&\\\hline
      Yellow & 1.333 & &&&\\\hline
      Green & 1.335 & &&&\\\hline
      Blue & 1.337 & &&&\\\hline
      Indigo & 1.340 & &&&\\\hline
      Violet & 1.344 & &&&\\\hline
    \end{tabular}
  \end{table}
\end{embeddedproblem}

A couple of comments about the rainbow problem are in order.

First, notice that although Calculus was crucial to the solution it wasn't
the \emph{whole} solution. Indeed, when we look back it seems that we
hardly used Calculus at all. We found the derivatives of $D(\alpha)$
and  $R(\alpha)$ and we used Newton's Method to find the roots of those
derivatives. And that's it. That's all of the Calculus we used.

Calculus was crucial, to be sure, but it is only one tool in our
toolbox. This will generally be true of the problems you encounter in
the future. There is a tendency to want to do ``Calculus'' because
this is a ``Calculus'' course. But this is a \underline{Mathematics}
course, first and foremost. You will need all of the mathematical
tools you have. 

% Second, in our comments on Proficiency~Drill~\ref{PD:DerivUnd} we
% observed that  Fermat's Theorem specifically requires that a function
% be differentiable at the point in question. This may seem like a
% pointless formality. Obviously, if the derivative doesn't exist at
% $x=a$ then we can't conclude that it is equal to zero at $x=a$. It
% hardly seems worth the effort of writing it down.

% However, this question of differentiability is more important than it
% appears to be right now. The fact is that not all functions have
% derivatives, and some functions have tangent lines but are not
% differentiable. We have glossed over this truth so far because we were
% learning the basics, but we are about to move beyond the basics so it is
% time to begin taking explicit notice of when a derivative exists, and
% when it does not. As a first approximation we will use this
% definition:

% \begin{mydefinition}{Non-Differentiability, First Approximation}
%   The derivative of $y=f(x)$ does not exist for any value of $x$ for
%   which $\dfdx{y}{x}$ is undefined.
%   \comment{As a practical matter this means
%   any value of $x$ that would result in dividing by zero.}
% \end{mydefinition}

% As we said this is a first approximation. This is a more subtle
% problem than it appears to be at first. We will refine this definition
% later.
\begin{embeddedproblem}
  Solve the following optimization problems. (That is find the maximum
  and minimum of the objective function, subject to the constraint.)
  Without context, you need to be mindful of where $\dfdx{f}{x}$ does
  and does not exist as a location for a possible extremum.
  \begin{enumerate}[label={\bf{}(\alph*)}]
    \begin{multicols}{2}
    \item Objective function: $f=2x+y $\\
      Constraint: $x^2+2y^2=2 $
    \item Objective function: $f=xe^y $\\
      Constraint: $x^2+y^2=2 $
    \item Objective function: $f=x+y $\\
      Constraint: $xy=4 $
    \item Objective function: $f=xy $\\
      Constraint: $x^2+9y^2=18 $
    \end{multicols}
  \end{enumerate}
\end{embeddedproblem}

% Ultimately, we will need conditions under which we can
% guarantee that the derivative exists, but for now we will simply take
% explicit notice  that sometimes it doesn't by stating differentiabilty
% as a requirement, when it is needed.



\begin{wrapfigure}[]{r}{2in}
\captionsetup{labelformat=empty}
\centerline{\includegraphics*[height=.9in,width=2in]{../Figures/Refraction11}}
\label{fig:Refraction11}
\end{wrapfigure}
Finally, when we graphed $F(\alpha)$ above we only graphed the part
where $0\le\alpha\le\frac{\pi}{2}$ because those are the only angles
that make sense in this problem. However if we were to graph it for a
wider range of values of $\alpha$ we would have a graph like the one
at the right. Notice that this graph has a maximum as well as a
minimum, but our procedure only tells us if we have a horizontal
tangent line. It does not distinguish between a horizontal tangent at a
minimum from a horizontal tangent at a maximum. If we had not known,
\emph{a priori}, that the minimum was between $0$ and $\frac{\pi}{2}$
how would we have known if we'd somehow found the maximum rather than
the minimum.

This seems like a minor problem at first, but it can become quite
problematic in more complex, or more abstract problems where we do not
have as much physical or geometric intuition to guide us as we had in
the rainbow problem. We will need some way to distinguish between a
maximum and a minimum analytically, without relying on pictures or our
intuition.


% \includegraphics*[height=3in,width=5in]{../Figures/DoubleRainbow1}

% \includegraphics*[height=3in,width=5in]{../Figures/DoubleRainbow0}



%%% Local Variables: 
%%% mode: latex
%%% outline-minor-mode: t
%%% eval:(flyspell-mode)
%%% TeX-master: "DiffCalc"
%%% End: 
